% !TeX program = pdflatex
%
% "Why Apery is alone: a scan of the modular irrationality race"
% Expository paper of the 2026-08-05 arc of River's odd-zeta programme.
% Primary sources: work/SCANNER_LEDGER.md (all eight sections),
% work/D1_CUSPIDAL_APPARATUS.md sec. 9, work/CATALAN_CONTROL.md,
% work/KZ_ELLIPTIC_REPORT.md, script work/z5eps/eps65_scanner.py.
%
% Evidence discipline: no claim is asserted here at a higher grade than in
% its source ledger. Finite computation is NEVER stated as proof.
%
\documentclass[11pt,reqno]{amsart}

\usepackage[T1]{fontenc}
\usepackage[utf8]{inputenc}
\usepackage{amsmath,amssymb,amsthm}
\usepackage{mathtools}
\usepackage{booktabs}
\usepackage{array}
\usepackage{enumitem}
\usepackage[margin=1.05in]{geometry}
\usepackage{microtype}
\usepackage[hidelinks]{hyperref}

\theoremstyle{plain}
\newtheorem{theorem}{Theorem}[section]
\newtheorem{proposition}[theorem]{Proposition}
\newtheorem{lemma}[theorem]{Lemma}
\newtheorem{corollary}[theorem]{Corollary}
\newtheorem{conjecture}[theorem]{Conjecture}

\theoremstyle{definition}
\newtheorem{definition}[theorem]{Definition}
\newtheorem{problem}[theorem]{Problem}

\theoremstyle{remark}
\newtheorem{remark}[theorem]{Remark}
\newtheorem{observation}[theorem]{Observation}
\newtheorem{lesson}[theorem]{Lesson}
\newtheorem{finding}[theorem]{Finding}

\newcommand{\Z}{\mathbb{Z}}
\newcommand{\Q}{\mathbb{Q}}
\newcommand{\R}{\mathbb{R}}
\newcommand{\C}{\mathbb{C}}
\newcommand{\HHH}{\mathbb{H}}
\newcommand{\eps}{\varepsilon}
\newcommand{\dn}{d_n}
\newcommand{\tc}{t_c}
\newcommand{\tcp}{t_c'}
\newcommand{\Ga}{\Gamma_0}

% evidence-class markers
\newcommand{\Thm}{\textup{\textsf{[THM]}}}
\newcommand{\Proved}{\textup{\textsf{[PROVED]}}}
\newcommand{\Verified}{\textup{\textsf{[VERIFIED]}}}
\newcommand{\VerifiedR}[1]{\textup{\textsf{[VERIFIED: #1]}}}
\newcommand{\Excluded}{\textup{\textsf{[EXCLUDED]}}}
\newcommand{\ExcludedR}[1]{\textup{\textsf{[EXCLUDED: #1]}}}
\newcommand{\Open}{\textup{\textsf{[OPEN]}}}

\begin{document}

\title[Why Ap\'ery is alone]{Why Ap\'ery is alone:\\ a scan of the modular irrationality race}

\author{[author]}

\date{August 5, 2026}

\begin{abstract}
Ap\'ery's 1978 proof that $\zeta(3)\notin\Q$ is usually described as a
miracle: a recurrence appears, its two solutions have the right
denominators, and a race between denominator growth and approximation
decay is won by a hair. We recall the race in a form that makes its two
competing exponents visible, and then --- following a scanner proposal of
our collaborator Sol --- we run the race over an entire class of
constructions: the companions attached to genus-zero modular curves with
an Atkin--Lehner fold.

Two things come out. First, a unification: the classical Ap\'ery
$\zeta(3)$ source itself \emph{vanishes at the fold},
$\Phi_\gamma(\tau_c)=0$ \VerifiedR{$10^{-122}$}, because it is the
$W_6$-odd Eisenstein direction. Ap\'ery's construction is therefore an
instance of exactly the mechanism used to manufacture cuspidal
companions --- an odd source vanishing at an Atkin--Lehner fixed point ---
realised in the Eisenstein class at a level whose fold ratio happens to
exceed $e^3$. Second, a bounded negative: over standard integral
eta-quotient hauptmoduln of level $N\le 25$, level $6$ is the
\emph{unique} winner of the $r=3$ race, and its odd cusp space
$S_4^-(\Ga(6))$ is zero; every level carrying odd cusp forms loses the
race by a factor $\ge 4$. The boundary levels $10$ and $14$, which sit at
score exactly $e^3$, are settled by direct computation as well. Along the
way the scan produces new arithmetic: the odd-cusp companion at level
$10$ has denominator exponent $\kappa=3$ at generic primes, with
deflation only at level primes, and an exact $5$-adic descent law
$v_5(B_{5m})=v_5(B_m)-2$ \VerifiedR{$m\le11$}. We state what is proved,
what is only computed, and the two problems that survive.
\end{abstract}

\maketitle

\tableofcontents

\section{Introduction: the race}
\label{sec:intro}

\subsection{Ap\'ery's theorem, in the shape we need}

In 1978 Ap\'ery announced that $\zeta(3)$ is irrational. The proof, as
digested by van der Poorten, Cohen and Beukers, fits on one page once one
knows the two sequences. Put
\[
  A_n \;=\; \sum_{k=0}^n \binom{n}{k}^2\binom{n+k}{k}^2,
\]
the \emph{Ap\'ery numbers} $1,5,73,1445,\dots$, and let $B_n$ be the
second solution of the same three-term recurrence
\begin{equation}
\label{eq:aperyrec}
  n^3 u_n \;=\; (34n^3-51n^2+27n-5)\,u_{n-1}\;-\;(n-1)^3u_{n-2},
\end{equation}
normalised by $B_0=0$, $B_1=6$. Then $A_n\in\Z$, the rational numbers
$B_n$ have denominators dividing $2\dn^3$ where
$\dn=\operatorname{lcm}(1,\dots,n)$, and
\[
  \frac{B_n}{A_n}\;\longrightarrow\;\zeta(3).
\]
So far this is only a fast approximation scheme. The proof is the
bookkeeping that follows. Three exponential rates govern everything:

\begin{itemize}[leftmargin=1.4em]
\item the \emph{growth rate} $\alpha$: $A_n\sim \alpha^n$ up to
  subexponential factors, here $\alpha=(1+\sqrt2)^4\approx 33.97$;
\item the \emph{decay rate} $\beta$: $|A_n\zeta(3)-B_n|\sim\beta^n$, here
  $\beta=(\sqrt2-1)^4\approx 0.0294=1/\alpha$;
\item the \emph{denominator rate} $\delta$: a common denominator for
  $A_n,B_n$ grows like $\delta^n$, here $\delta=e^3$ because
  $\dn\sim e^n$ by the prime number theorem and the exponent is $3$.
\end{itemize}

Clearing denominators produces genuine integers $\tilde A_n=\dn^3A_n$ and
$\tilde B_n = 2\dn^3B_n$ with
\[
  0\;<\;|\tilde A_n\zeta(3)-\tilde B_n|\;\approx\;(\delta\beta)^n .
\]
If $\zeta(3)=p/q$ were rational, the left-hand side would be a nonzero
integer multiple of $1/q$, hence bounded below by $1/q$; but the
right-hand side tends to $0$ as soon as $\delta\beta<1$. Contradiction.
Everything therefore rests on a single inequality.

\begin{definition}[The race condition; Sol's score]
\label{def:score}
For a construction with rates $(\alpha,\beta,\delta)$ as above, the
\emph{score} is
\[
  \mathrm{score} \;:=\; \delta\,|\beta| ,
\]
and the construction \emph{wins the race} (i.e.\ proves irrationality of
its limit) precisely when $\mathrm{score}<1$.
\end{definition}

For Ap\'ery's $\zeta(3)$: $\delta|\beta| = e^3\cdot 0.0294 \approx
20.09/33.97 \approx 0.59 <1$. The margin is $41\%$; there is no slack to
spare and none of the three rates could be much worse. For Ap\'ery's
$\zeta(2)$ construction (the $\binom{n}{k}^2\binom{n+k}{k}$ family) the
corresponding numbers are $\delta=e^2$, $|\beta|^{-1}=\varphi^{10}
\approx 11.09$, and the score is $0.67$.

\begin{remark}
The formulation of the race as a one-number \emph{score} attached to a
construction --- and the proposal to scan it systematically over a class of
constructions rather than admire it in the two classical cases --- is due
to our collaborator Sol. The entire paper is an execution of that
proposal. We keep the terminology.
\end{remark}

\subsection{Why the score is computable in the modular world}

The reason a scan is possible at all is that for the constructions we
care about --- the ones attached to modular curves --- the three rates are
not independent mysteries. They are read off from one picture.

Suppose $\Gamma$ is a genus-zero congruence group with hauptmodul $t$
(a generator of its function field, normalised to $t=q+O(q^2)$), and
suppose $F$ is a weight-$k$ modular form on $\Gamma$ that, written as a
function of $t$, satisfies a Picard--Fuchs equation
$L\,F=0$ whose coefficient sequence $A_n=[t^n]F$ is integral. In the
sporadic cases this is exactly the classical dictionary between Ap\'ery-like
sequences and modular parametrisations. The differential operator $L$ has
singularities at the roots of a polynomial $P(t)$; write $\tc$ for the
root nearest the origin (the \emph{fold}) and $\tcp$ for the conjugate
root. Then:

\begin{itemize}[leftmargin=1.4em]
\item $\alpha = 1/|\tc|$ --- the growth of $A_n$ is governed by the
  nearest singularity;
\item if the companion $B_n$ is built from a \emph{source} that vanishes
  at the fold, the particular solution's own singularity is pushed out to
  $\tcp$, so the error decays like $|\tcp|^{-n}$, i.e.\
  $|\beta| = 1/|\tcp|$;
\item $\delta=e^r$ where $r$ is the denominator exponent: $r=3$ for a
  weight-$4$ source over a weight-$2$ family, $r=2$ for weight-$3$ over
  weight-$1$.
\end{itemize}

Hence

\begin{equation}
\label{eq:score}
  \boxed{\;\mathrm{score}_r \;=\; \frac{e^r}{|\tcp|}\;}
\end{equation}

and \emph{the race is won exactly when the conjugate singularity of the
Picard--Fuchs operator is farther from the origin than $e^r$}. That is a
statement about one complex number attached to a level, and one can
enumerate levels.

For Ap\'ery's $\zeta(3)$ the relevant curve is level $6$ with
$P(t)=1-34t+t^2$, whose roots are $\tc=17-12\sqrt2\approx 0.0294$ and
$\tcp = 17+12\sqrt2\approx 33.97$, giving score $20.09/33.97=0.59$ as
above. The whole of \S\ref{sec:scan} is the systematic version of this
one line.

\subsection{What this paper contains}

\S\ref{sec:unification} explains the structural discovery that makes the
scan meaningful: Ap\'ery's source vanishes at the fold, so his
construction is the \emph{same} construction as the deliberately
manufactured cuspidal apparatus of the companion paper --- one mechanism
in two representation-theoretic classes. \S\ref{sec:scan} runs the scan
and reports the verdict, together with a lesson about normalisation
dependence that changes what the optimisation problem even is.
\S\ref{sec:s1} settles the two boundary levels. \S\ref{sec:s2} reports
the arithmetic that the scan turned up as a by-product --- a prime-local
denominator law and an exact $5$-adic descent law --- and closes the
``arithmetic loophole''. \S\ref{sec:limits} states honestly what is
bounded, assumed, and unproved.

The intended reader is a graduate student who has seen Ap\'ery's proof
once and knows what a modular form is. We have kept the modular
prerequisites to: hauptmoduln, Atkin--Lehner involutions, oldforms, and
the dimension formula for $S_k(\Ga(N))$.

\subsection{Evidence discipline}

This programme labels every assertion. \Proved{} means a paper proof
exists; \Thm{} means proved and adversarially refereed; \VerifiedR{range}
means exact computation over a stated finite range and is \emph{never}
treated as proof; \Excluded{} means a proved negative; \Open{} means open,
stated in its sharpest known form. Nothing in this paper is asserted at a
higher grade than in its source ledger. In particular the central
negative result of \S\ref{sec:scan} is a bounded computational verdict
over an explicitly delimited search domain, not a theorem about all
modular constructions --- and we say so again in \S\ref{sec:limits}.

\section{The unification: Ap\'ery's source vanishes at the fold}
\label{sec:unification}

\subsection{Sources and companions}

Fix a genus-zero level $N$, a hauptmodul $t$, and a weight-$2$ modular
form $F$ with $A_n=[t^n]F\in\Z$ and Picard--Fuchs operator $L$. A
\emph{companion} of $F$ is a second solution: a sequence $B_n$ satisfying
the same recurrence with an inhomogeneity. The standard way to produce
one is to pick a modular form $g$ of weight $r+1$ (the \emph{source}) and
set
\[
  B \;=\; F\,\theta_q^{-r} g, \qquad \theta_q = q\frac{d}{dq},
\]
where $\theta_q^{-r}$ is the $r$-fold formal antiderivative in $q$
(an Eichler integral). Expanding $B$ in powers of $t$ gives $B_n$, and
the operator applied to $B$ returns a computable inhomogeneity: in the
normalised form used throughout this programme,
\[
  L\bigl(F\theta_q^{-r}g\bigr) \;=\; \Phi\quad\text{with}\quad
  \Phi \;=\; \frac{t\,\sigma^r}{P\,F}\cdot(\dots),
\]
where $\sigma=\theta_q t/t$. The point of the modular-anchor programme is
that $\Phi$ is again a modular object, and that for all twelve
identifiable sporadic families it is a \emph{pure Eisenstein series} of
weight $r+1$ \VerifiedR{$q^{60}$, held-out band $q^{27}$--$q^{60}$}.

Ap\'ery's $B_n$ is the case $N=6$, $r=3$, with source $\Phi_\gamma$ the
weight-$4$ Eisenstein direction on $\Ga(6)$.

\subsection{The vanishing mechanism, first seen in the cuspidal case}

The mechanism was isolated in the reverse direction. On Domb's curve
(level $12$, $\dim S_4(\Ga(12))=2$, spanned by the two embeddings
$f_6(q), f_6(q^2)$ of the level-$6$ newform $f_6=(\eta_1\eta_2\eta_3\eta_6)^2$),
one has the Atkin--Lehner action
\[
  (f_6|_4W_{12})(\tau) \;=\; 4f_6(2\tau) \qquad \Proved,
\]
which makes
\[
  f^{*} \;=\; f_6(q)-4f_6(q^2)\;=\;q-6q^2-3q^3+12q^4+6q^5+\cdots
  \;\in\;S_4(\Ga(12))
\]
the $W_{12}$-\emph{odd} eigenvector. Because the weight-$4$ automorphy
factor at the fixed point $\tau_{12}=i/\sqrt{12}$ equals $1$, oddness
forces $f^{*}(\tau_{12}) = -f^{*}(\tau_{12})$, i.e.
\[
  f^{*}(\tau_{12})\;=\;0 .
\]
That single vanishing does two independent jobs (see the companion paper
on the cuspidal apparatus, and \cite[\S4]{D1}):

\begin{enumerate}[label=(\roman*),leftmargin=2em]
\item \textbf{Exponential convergence.} The inhomogeneity
  $t f^{*}/\Phi$ becomes regular at the fold, so the particular solution
  is singular only at the \emph{conjugate} root; the error ratio is
  $(\tc/\tcp)^n$ instead of $O(1)$.
\item \textbf{Clean limit.} For an $\eps=-1$ source, the fixed-point
  Eichler equation selects value-plus-quasiderivative data that the odd
  period polynomial reduces to a single critical $L$-value; the
  second-kind (quasiperiod) term is structurally absent.
\end{enumerate}

The resulting apparatus has $\dn^3B^{*}_n\in\Z$ \VerifiedR{$n\le40$},
error $\approx 4^{-n}$, and
$\lim B^{*}_n/A_n = L(f_6,3)/2$ \VerifiedR{PSLQ, 61 digits}. Its score,
however, is $e^3/4\approx 5.0>1$: it does \emph{not} prove irrationality
of $L(f_6,3)$. It was built to test whether cuspidal $L$-values could be
brought inside the machinery at all. They can.

\subsection{The discovery}

Having isolated ``odd source vanishing at the Atkin--Lehner fold'' as the
mechanism, one naturally asks what Ap\'ery's own source does at his fold.
The answer:

\begin{finding}[$\Phi_\gamma(\tau_c)=0$]
\label{find:phi}
For the level-$6$ family $\gamma$ (Ap\'ery's $\zeta(3)$ construction), the
source $\Phi_\gamma$ vanishes at the Atkin--Lehner fixed point
$\tau_c = i/\sqrt6$: \[\Phi_\gamma(\tau_c)=0\qquad
\VerifiedR{$|\Phi_\gamma(\tau_c)|<10^{-122}$}.\]
The reason is the same parity argument as for $f^{*}$: $\Phi_\gamma$ lies
in the $W_6$-\emph{odd} part of the weight-$4$ space, and any $W_N$-odd
weight-$4$ form vanishes at the fixed point because the automorphy factor
there is $+1$.
\end{finding}

We label the numerical statement \Verified{} and the parity explanation
\Proved{} (it is the two-line argument just given, once the eigenvalue
$\eps(\Phi_\gamma,W_6)=-1$ is known; the eigenvalue itself is read off
the Eisenstein decomposition).

The consequence is conceptual, and it is the reason this paper exists:

\begin{observation}[One mechanism, two classes]
\label{obs:unify}
Ap\'ery's $\zeta(3)$ construction is \emph{not} a separate miracle from
the manufactured cuspidal apparatus. Both are instances of:
\begin{center}
\emph{an odd source, vanishing at an Atkin--Lehner-fixed fold, integrated
against a modular family.}
\end{center}
They differ only in which part of the weight-$(r+1)$ space the source
lives in --- Eisenstein or cuspidal --- and in whether the level's fold
ratio beats $e^r$.
\end{observation}

This organises the landscape into a $2\times2$ table.

\begin{table}[h]
\centering
\renewcommand{\arraystretch}{1.35}
\begin{tabular}{@{}l>{\raggedright\arraybackslash}p{4.6cm}>{\raggedright\arraybackslash}p{5.0cm}@{}}
\toprule
 & \textbf{odd Eisenstein source} & \textbf{odd cusp source} \\
\midrule
\textbf{level 6} \newline ($|\tcp|=33.97>e^3$) &
  \textbf{Ap\'ery's $\zeta(3)$ proof} --- score $0.59$, wins &
  none exists: $S_4^-(\Ga(6))=0$ \newline ($f_6$ is Fricke-\emph{even}) \\
\textbf{level 12} \newline ($|\tcp|=4<e^3$) &
  a $\zeta(3)$-type companion; loses the race &
  the $L(f_6,3)/2$ apparatus; loses the race (score $\approx5.0$) \\
\bottomrule
\end{tabular}
\medskip
\caption{The mechanism in its two classes at the two relevant levels.}
\label{tab:2x2}
\end{table}

Reading Table~\ref{tab:2x2} the obvious question writes itself.

\begin{problem}[The cuspidal-Ap\'ery cell]
\label{prob:cell}
Is there a genus-zero group $\Gamma$, with a rectified family whose
Atkin--Lehner-fixed fold has $|\tcp|>e^r$, and with
$S^-_{r+1}(\Gamma)\ne 0$? Such a level would give a race-winning
apparatus whose limit is a critical $L$-value of a cusp form --- an
irrationality proof for a cuspidal period.
\end{problem}

Both diagonal cells of Table~\ref{tab:2x2} are occupied; the top-right
cell is empty for a reason specific to level $6$ ($\dim S_4=1$ and that
one form has the wrong parity). Is the emptiness an accident of small
level? \S\ref{sec:scan} answers.

\begin{remark}[Two repairs that do not work]
Two obvious fixes are excluded before the scan starts. (i) At level $6$
one cannot subtract an Eisenstein multiple to force a cusp form to vanish
at the fold, because $\Phi_\gamma$ already vanishes there --- the
correction term one would want is identically zero on the fold
\Verified. (ii) Moving to a level with a complex fold does not help: the
regularisation argument requires a \emph{real} fold at an Atkin--Lehner
fixed point, as the $\delta$-family control test showed (there the fold
is at $q_c\approx 0.1137-0.1970i$, $f^{*}$ does not vanish, and the
companion's coefficient ratios oscillate with both conjugate folds
singular).
\end{remark}

\subsection{Why the modular class is the right arena}

A reader might object that we are scanning a small pond. Two control
experiments justify the choice of arena; both are negative results about
the \emph{hypergeometric} alternative.

\emph{Zudilin's Catalan recurrence} (a well-poised ${}_6F_5$-world
construction with limit Catalan's constant $G$) satisfies the universal
boundary-defect structure of our formalism --- $L(y_v)=\tfrac{13}{2}t$
exactly \VerifiedR{$q^{40}$} --- but fails both halves of the
modular-anchor mechanism: its canonical nome admits \emph{no} integralising
rational rescale (exhaustive scan over $\mu=\pm2^a3^b5^c$, $|a|\le8$,
$|b|\le3$, $|c|\le2$), and its fold connection value misses $(2/13)G$ by
$3.2\times10^{-2}$, far above truncation error. The diagnosis is
structural: the raw operator has order six with half-integer indicial
pairs $\{0,0,\tfrac12,\tfrac12,\cdot,\cdot\}$. \ExcludedR{bounded}

\emph{Koutschan--Zudilin's elliptic $L$-value construction} --- the unique
one in the literature realising elliptic $L$-values as Ap\'ery limits ---
carries the same half-integer structure ($x^{n-1/2}$ integrands,
$(2n\pm1)$ coefficients). Executed at $z=1/2$ (the conductor-$32$ CM
case, where CM might have saved it), modular rectification fails
maximally: canonical-nome denominators $1352, 1.7\times10^{13},
1.9\times10^{25}$ at $n=2,3,4$, no integralising rescale, and the
indicial polynomial $\theta(\theta+1)^2(2\theta-1)^2(2\theta+1)^2\cdots$
has exponent $0$ only \emph{simple}, so no Frobenius log-partner exists
and the canonical coordinate is not even well-founded. \ExcludedR{bounded}

\begin{lesson}
The source of a companion depends on the \emph{geometric realisation},
not on the constant. The same $G$ is reached inside the modular class
(sporadic family E, level $8$, pure weight-$3$ $\chi_{-4}$-Eisenstein
source, integral mirror map) and outside it (Zudilin's recurrence, no
modular structure in the canonical coordinate). Denominator and
congruence phenomena belong to the realisation.
\end{lesson}

Consequently: the known non-modular $L$-value recurrences are
hypergeometric-realised and outside our mechanism, while the mechanism
itself has exactly one known race winner per weight. Scanning the modular
class is scanning the place where the winners live.

\section{The scan}
\label{sec:scan}

\subsection{Method}

The scanner (\texttt{work/z5eps/eps65\_scanner.py}) implements
\eqref{eq:score} directly. For each level $N$ in a list of levels
carrying a standard integral eta-quotient hauptmodul $t$, it computes:

\begin{enumerate}[label=\textbf{Step \arabic*.},leftmargin=3.4em]
\item \textbf{The hauptmodul.} $t=\prod_m\eta(m\tau)^{e_m}$ with the
  standard exponent vectors --- e.g.\ level $6$:
  $(\eta_1\eta_6/\eta_2\eta_3)^{12}$; level $12$:
  $(\eta_1\eta_{12}/\eta_3\eta_4)^{4}$; level $10$:
  $(\eta_1\eta_{10}/\eta_2\eta_5)^{6}$; prime levels
  $(\eta_1/\eta_N)^{24/(N-1)}$.
\item \textbf{The Atkin--Lehner M\"obius map.} $t\circ W_N$ is again a
  hauptmodul, hence $t\circ W_N=(at+b)/(ct+d)$ for constants determined
  by the level. These are fitted numerically from samples
  $\tau=iy$, $y\in\{0.35,0.45,0.55,0.65\}$, with the fourth sample used
  as a consistency check (residuals reported per level).
\item \textbf{The critical values.} The singular values of the operator
  are the critical values of $t$, i.e.\ the values of $t$ at the zeros of
  \[
    \sigma/t \;=\; \tfrac{1}{24}\sum_m e_m\,m\,E_2(q^m),
  \]
  found by seeded Newton iteration over the complex disc --- both the real
  folds and the complex second classes. Equivalently, and as a
  cross-check, they are the fixed points of the M\"obius involution of
  Step 2: solving $ct^2+(d-a)t-b=0$ gives $\{\tc,\tcp\}$, and $\tc$ is
  identified as the one nearest $t(i/\sqrt N)$.
\item \textbf{The score and the cusp dimension.} Output
  $\mathrm{score}_r=e^r/|\tcp|$ for $r=2,3$, together with
  $\dim S_4(\Ga(N))$ from the classical genus/elliptic-point formula.
\end{enumerate}

\subsection{Validation anchors}

A scan is worth only its anchors. The following were reproduced
\emph{exactly} by the pipeline before any new level was trusted:

\begin{itemize}[leftmargin=1.4em]
\item \textbf{Level 6} (Ap\'ery's $\zeta(3)$): $\tc=0.0294\ldots$,
  $|\tcp|=33.97\ldots$ --- the roots of $1-34t+t^2$, known from the family
  operator.
\item \textbf{Level 12} (Domb): $\tc=1/16$, $|\tcp|=1/4$ --- the roots of
  $P_\alpha=(1-4t)(1-16t)$.
\item \textbf{Level 8} (family $\eps$): $(0.0429\ldots,\,1.457\ldots)$.
\item \textbf{Level 20} (family $\eta$): the complex fold
  $0.088\pm0.016i$, matching the known complex-fold classification (disc
  $-16$).
\item \textbf{Level 5}: the symmetric $\pm1/\sqrt5$-scale pair, i.e.\
  $|\tc|=|\tcp|$, consistent with $t\mapsto\kappa/t$.
\end{itemize}

\subsection{The table}

Table~\ref{tab:scan} is the full first run, at $r=3$
($\mathrm{score}_3=e^3/|\tcp|\approx 20.09/|\tcp|$; win iff $<1$).

\begin{table}[h]
\centering
\renewcommand{\arraystretch}{1.2}
\begin{tabular}{@{}llllll@{}}
\toprule
$N$ & $\dim S_4$ & $\tc$ & $|\tcp|$ & $\mathrm{score}_3$ & verdict\\
\midrule
$2,3,4,5,$ & $0$--$5$ & $\pm\sqrt\kappa$ & $=|\tc|$ & --- & no analytic gain\\
$7,9,13,25$ & & symmetric & & & ($t\mapsto\kappa/t$)\\
$6$ & $1$ & $0.0294$ & $\mathbf{33.97}$ & $\mathbf{0.59}$ &
   \textbf{WIN} --- but $S_4^-(\Ga(6))=0$\\
$8$ & $1$ & $0.0429$ & $1.457$ & $13.8$ & fail\\
$10$ & $3$ & $0.0557$ & $1.000$ & $20.09$ & fail (exactly $e^3$)\\
$12$ & $3$ & $0.0625$ & $0.25$ & $80$ & fail (Domb; the $L(f_6,3)/2$ apparatus)\\
$14$ & $4$ & $0.1497$ & $1.000$ & $20.09$ & fail (exactly $e^3$)\\
$15$ & $4$ & $0.2011$ & $0.726$ & $27.7$ & fail\\
$18$ & $5$ & $0.0895$ & $\lesssim1$ & $\gg1$ & fail (second class not fully resolved)\\
$20$ & $6$ & complex fold & --- & --- & excluded (no real AL fold)\\
$24$ & $8$ & $0.0981$ & $0.545$ & $36.8$ & fail\\
\bottomrule
\end{tabular}
\medskip
\caption{The $r=3$ scan over standard integral eta hauptmoduln,
$N\le25$. \VerifiedR{numeric, $50$ digits working precision}}
\label{tab:scan}
\end{table}

The $r=2$ scan (score $e^2/|\tcp|$, relevant for weight-$3$ sources over
weight-$1$ families) produced no winner either, subject to the
normalisation caveat of \S\ref{sec:normalisation}.

\subsection{The verdict}

\begin{finding}[The cuspidal-Ap\'ery cell is empty in the scanned domain]
\label{find:verdict}
Within the natural search domain --- standard integral eta hauptmoduln,
levels $N\le25$ --- Problem~\ref{prob:cell} has no solution.
Level $6$ is the unique $r=3$ race winner, and its odd cusp space
vanishes; every level carrying odd cusp forms loses the race by a factor
$\ge 4$. \VerifiedR{Table~\ref{tab:scan}; domain as stated}
\end{finding}

Or, in one sentence: \emph{Ap\'ery's $\zeta(3)$ apparatus is the unique
odd-vanishing race-winning construction in its class --- and it is
Eisenstein because it has to be.} The classical picture, in which the
Eisenstein-ness of Ap\'ery's source looks like a lucky feature, inverts:
the cuspidal alternative is not available at any level where the race can
be won.

The reader should note carefully what kind of statement
Finding~\ref{find:verdict} is. It is a verdict over an enumerated finite
domain, computed numerically. It is \emph{not} a theorem that no
race-winning cuspidal construction exists. \S\ref{sec:limits} lists
precisely which doors it leaves open, and \S\ref{sec:normalisation}
explains why one of those doors is much wider than it first appears.

\subsection{Two boundary curiosities}

Levels $10$ and $14$ sit at $|\tcp|=1.000$ to all computed digits ---
their conjugate class lies on the unit $t$-circle --- hence at score
\emph{exactly} $e^3$. This is the knife edge: the analysis contributes
nothing either way, and the outcome is decided entirely by arithmetic.
Since $\dim S_4=3$ and $4$ there, odd cusp vectors could exist. Any
denominator improvement whatsoever --- a $\dn^{3-\eps}$ law of the sort the
sporadic ``deflation'' phenomena hint at --- would tip these levels into
the win column and produce a cuspidal irrationality proof. This was the
sharpest target the first run produced, recorded as:

\begin{problem}[S-1]
\label{prob:s1}
Determine the exact odd-cusp dimensions at levels $10$ and $14$, and
whether any integral normalisation or auxiliary-source choice achieves
denominator growth strictly below $e^{3n}$.
\end{problem}

\S\ref{sec:s1} answers it.

\subsection{The normalisation lesson}
\label{sec:normalisation}

Before that, a lesson that reshapes the problem. Look again at the first
row of Table~\ref{tab:scan}: at prime level the standard hauptmodul is
\emph{symmetric}, $t\circ W_N=\kappa/t$, so its two fixed points are
$\pm\sqrt\kappa$ and $|\tc|=|\tcp|$ identically. There is zero analytic
gain \emph{by construction}. Yet the classical Ap\'ery $\zeta(2)$ proof
lives at level $5$ and does win its race.

The resolution: the $\zeta(2)$ win exists in the $\Gamma_1(5)$
(asymmetric) normalisation of the level-$5$ hauptmodul, where
$|\tcp|/|\tc| = \varphi^{10}\approx123$; in the symmetric eta
normalisation the same level shows no gain at all. Both are integral
coordinates on the same curve.

\begin{lesson}[Normalisation non-invariance]
\label{lesson:norm}
The race score is \emph{not} an invariant of the level. It is an
invariant of the \emph{integral normalisation of the hauptmodul}. Two
integral coordinates on the same modular curve can give scores on
opposite sides of $1$.
\end{lesson}

This is a genuine two-coordinate comparison at a single level, so it
settles the question of principle: score is coordinate-dependent. The
consequence for Problem~\ref{prob:cell} is that the true optimisation
domain is not the (small, enumerable) set of levels but the
\emph{lattice of integral M\"obius changes of $t$}, each trading
singularity positions against denominator growth --- one cannot move
$|\tcp|$ outward for free, because rescaling $t$ rescales $A_n$ and
$B_n$ and hence the denominators. This is precisely Sol's
``extremal function'' formulation, now with data attached. We return to it
as the surviving open problem in \S\ref{sec:s2}.

\section{S-1: the level-$10$ boundary test}
\label{sec:s1}

Problem~\ref{prob:s1} asks whether the knife-edge levels can be tipped.
We report the level-$10$ computation in full, because every layer of the
mechanism behaves \emph{exactly} as the theory predicts --- and the
construction still fails, for a reason the first framing of S-1 had
mis-stated.

\subsection{The odd vector exists}

The level-$10$ weight-$4$ oldspace contains the embeddings of
$f_5=\eta_1^4\eta_5^4$, the level-$5$ newform. Its Atkin--Lehner
eigenvalue is $\eps_5=+1$, read off from $a_5=-5$; the general oldspace
lemma (whose two previous instances are levels $12$ and $6$) gives
$c=\eps p^{k/2}$, here $c=4$, so
\[
  (f_5|_4W_{10}) = 4f_5(q^2),
\]
and therefore
\[
  \boxed{\,f^- \;=\; f_5(q)-4f_5(q^2)\,}
\]
is the $W_{10}$-odd eigenvector. It vanishes at the fold automatically,
by the same automorphy-factor parity argument as in
\S\ref{sec:unification}. This is the \emph{third} instance of the
oldspace vanishing lemma and the exact structural analogue of level
$12$'s $f^{*}$. \Verified

\subsection{The families are rectified}

Level $10$ supports integral principal solutions: for
$F=2E_2(q^2)-E_2(q)$ one gets $A_n = 1,24,168,1752,\dots\in\Z$, and for
$F=(5E_2(q^5)-E_2(q))/4$ one gets $1,6,54,582,\dots\in\Z$. Nothing about
level $10$ is degenerate. \VerifiedR{exact, series to $q^{36}$}

\subsection{Integrality holds}

The odd-cusp companion $B_n=[t^n]F\theta_q^{-3}f^-$ satisfies
\[
  \dn^3 B_n \in \Z \qquad \VerifiedR{$n\le36$, exact}.
\]
The $\dn^3$-integrality layer --- which the programme has now observed at
levels $6$, $12$ and $10$, cuspidally and Eisenstein-wise, on several
curves --- holds yet again. (It remains \emph{observed}, not proved, in
the cuspidal case; see \S\ref{sec:limits}.)

\subsection{And yet the construction fails}

Two ways, both fatal:

\begin{itemize}[leftmargin=1.4em]
\item \textbf{The denominators grow.} $\log(\mathrm{den}\,B_n)/n \approx
  2.27$--$2.43$ for $n=20$--$36$: slightly below the nominal $\dn^3$-rate
  $3$, but nowhere near bounded. There is no $\dn^{3-\eps}$ miracle with
  $\eps$ big enough to matter, and certainly no boundedness.
\item \textbf{The forms do not decay.} At $|\tcp|=1.000$ the linear forms
  $A_n\xi-B_n$ are $O(1)$: they do not tend to $0$ at all. A score of
  exactly $1$ does not mean ``arithmetic decides''; it means the analytic
  input is \emph{zero}, and one would need bounded denominators
  \emph{and} strictly decaying forms. Level $10$ has neither.
\end{itemize}

\begin{finding}[S-1 answered: NO]
\label{find:s1}
The cuspidal-Ap\'ery cell is empty in the full scanned domain, now
including the boundary levels. \VerifiedR{level $10$ exact to $q^{36}$;
level $14$ by the same unit-circle argument}
\end{finding}

For level $14$ we note a correction: its hauptmodul was mis-specified in
the first run; the correct one is $(\eta_2\eta_{14}/\eta_1\eta_7)^3$, and
its conjugate class sits at the same unit-circle scale, so the same
argument applies.

\begin{remark}[What survives of S-1]
Only the normalisation-optimisation question of
Lesson~\ref{lesson:norm}: whether some integral M\"obius change of
coordinate at \emph{any} level can push $|\tcp|$ above the observed
denominator rate. At the time this finding was recorded, the empirical
rate for odd-cusp companions looked like $e^{2.3n}$ rather than $e^{3n}$,
which would have lowered the bar --- a level with $|\tcp|>e^{2.3}\approx
10$ and odd cusp forms is not obviously excluded. That deflation was
real in the data and unexplained. Resolving it was posed as S-2, and it
is the subject of the next section.
\end{remark}

\section{S-2: the prime-local denominator law, and the closing of the loophole}
\label{sec:s2}

\subsection{The question}

Is the odd-cusp companion's true denominator exponent $\kappa=3$, or is
it genuinely smaller? Everything hinges on this: the threshold to beat is
$e^\kappa$, and the difference between $e^3\approx20.1$ and
$e^{2.3}\approx10.0$ is the difference between ``no candidate levels''
and ``several''.

The protocol: compute $B_n$ exactly for $n\le72$ for both integral
choices of $F$ at level $10$, and split the denominator prime by prime.
Write $h_p(n)=v_p(\dn)=\lfloor \log_p n\rfloor$ for the $p$-adic
valuation of $\operatorname{lcm}(1,\dots,n)$, and $e_p(n)=-v_p(B_n)$ for
the power of $p$ actually appearing in the denominator of $B_n$. The
nominal law is $e_p\le 3h_p$; the question is where it is sharp.

\subsection{The law}

\begin{finding}[Prime-local denominator law at level 10]
\label{find:kappa}
Uniformly across both integral $F$'s, for $n\le72$
\VerifiedR{exact rational arithmetic}:
\begin{align}
  e_p(n) &= 3h_p(n) &&\text{for } p\nmid 10 \quad\text{(generic primes:
   the \emph{full} $\dn^3$ cost)},\label{eq:generic}\\
  e_5(n) &= 2h_5(n) &&\text{exactly},\label{eq:five}\\
  e_2(n) &\ll 2h_2(n) &&\text{(the strongest deflation).}\label{eq:two}
\end{align}
Concretely for \eqref{eq:generic}: $e_7=3$ at $n=7$, $e_7=6$ at $n=49$,
and likewise for $p=11,13$. For \eqref{eq:five}: $e_5=2$ at $n=5$,
$e_5=4$ at $n=25$. The auxiliary quantity
$Q_n=\mathrm{den}(B_n)/\gcd(\mathrm{den}(B_n),\dn^2)$ contains every
generic prime $\le n$ to power exactly $1$ --- which is precisely the
difference between $3h_p$ and $2h_p$ --- and contains $5$ never.
\end{finding}

Two consequences, one negative and one positive.

\begin{corollary}[$\kappa=3$; the $2.3$ was an artifact]
The globally fitted exponent $2.27$ of \S\ref{sec:s1} was a finite-size
effect. It \emph{rises} with $n$ --- reaching $2.75$ by $n=72$ --- tracking
$3\psi(n)/n$ minus $O(\log n)$ worth of level-prime savings, where
$\psi$ is Chebyshev's function. The true exponent is
$\boxed{\kappa=3}$, and the analytic threshold stays $e^3$.
\end{corollary}

Hence the candidate universal ``$\dn^2$-type'' theorem --- which would have
made the boundary levels winnable --- is \Excluded{} at generic primes.
The scanner verdict of \S\ref{sec:scan} stands, and now stands
arithmetically as well as analytically: \emph{the cuspidal-Ap\'ery cell
remains empty, with the arithmetic loophole closed}.

\subsection{What is gained: level-prime deflation}

The positive consequence is new arithmetic. The deflations at $p\mid N$
are real, structured, and provable-looking --- they are the local behaviour
of Frobenius on the Eichler extension at the primes of bad reduction. The
$5$-adic one turns out to have an exceptionally clean form. Probing
$n\le60$ exactly:

\begin{finding}[The exact $5$-adic descent law]
\label{find:descent}
For the level-$10$ odd-cusp companion,
\[
  \boxed{\;v_5(B_{5m}) \;=\; v_5(B_m)-2\;}
\]
\VerifiedR{all tested $m\le11$, uniformly, including cases with
$v_5(B_m)<0$}.
That is: each $5$-adic digit of $n$ costs exactly \emph{two} powers of
$5$. The mechanism visible in the data: the division by $m^3$ in the
recurrence loses three powers, and the numerator of the forced recurrence
recovers exactly one.
\end{finding}

Everything else about the $5$-adic side follows from
Finding~\ref{find:descent} by digit bookkeeping: the equality set of
$e_5=2h_5$ (which holds for $34$ of the $60$ values $n\le60$) and its
exception windows --- the $5$-adic digit patterns
$21_5$--$24_5$, $41_5$--$44_5$, $111_5$--$114_5$, \dots --- are exactly
what the descent law predicts. Moreover $25\cdot B_{5m}\bmod 5$ is a unit
whose value tracks the pair $(A,B)$: this is the digit-transition
constant of a mixed Lucas congruence, one experiment away from explicit.

Note the shape of the law: it is an \emph{equality}, a $5$-adic
recursion, not an inequality. That is unusual and is what makes it look
provable.

\begin{problem}[S-3, open]
\label{prob:s3}
Prove $v_5(B_{5m})=v_5(B_m)-2$ for all $m$. Two routes are identified:
the source side (the action of $U_5$ on $\theta_q^{-3}f^-$, which would
explain why the law is independent of the choice of $F$ --- the observed
$F$-invariance is otherwise unexplained), or the forced recurrence
directly (shortest). \Open
\end{problem}

\subsection{The surviving extremal problem}

With the arithmetic loophole closed, exactly one route to a cuspidal
irrationality proof survives inside this mechanism, and it is the one
Lesson~\ref{lesson:norm} opened.

\begin{problem}[The integral-model extremal problem]
\label{prob:extremal}
\begin{center}
\fbox{\parbox{0.88\linewidth}{\medskip
Find a quadruple $(N,\,t,\,F,\,f^-)$ --- a level, an \emph{integral}
coordinate $t$ on the modular curve, a rectified weight-$2$ family $F$
with integral $t$-coefficients, and a $W_N$-odd cusp form $f^-$ of weight
$4$ --- such that the conjugate singular value satisfies
\[
  |\tcp| \;>\; e^{3}\approx 20.09
\]
in that coordinate, with the $\dn^3$-integrality of the companion
preserved. \Open\medskip}}
\end{center}
Equivalently: optimise the score over the lattice of integral M\"obius
changes of coordinate, where every outward push of $|\tcp|$ is paid for
in denominator growth.
\end{problem}

This is a genuine extremal problem, not a restatement of
Problem~\ref{prob:cell}: the scan fixed one standard coordinate per
level, and Lesson~\ref{lesson:norm} proves that this choice can decide
the race. Level $5$ is the existence proof that the coordinate lattice
matters --- there the asymmetric normalisation gains a factor
$\varphi^{10}\approx123$ over the symmetric one.

Together with Problem~\ref{prob:s3}, this is what the arc leaves open.
Everything else in the cell is settled, for the scanned class:
analytically by Table~\ref{tab:scan}, at the boundary by
Finding~\ref{find:s1}, and arithmetically by Finding~\ref{find:kappa}.

\subsection{Summary of the position}

\begin{table}[h]
\centering
\renewcommand{\arraystretch}{1.25}
\begin{tabular}{@{}p{6.6cm}l@{}}
\toprule
statement & grade\\
\midrule
$\Phi_\gamma(\tau_c)=0$ (Ap\'ery's source is odd, vanishes at the fold)
  & \VerifiedR{$10^{-122}$} + \Proved{} parity argument\\
$(f|_kW_{Np})=\eps p^{k/2}f(q^p)$ on the oldspace (three instances)
  & \Proved{} (level 12), \Verified{} (levels 6, 10)\\
Level 6 is the unique $r=3$ winner, $N\le25$, standard hauptmoduln
  & \VerifiedR{Table~\ref{tab:scan}}\\
$S_4^-(\Ga(6))=0$
  & \Proved{} ($\dim S_4=1$, $f_6$ Fricke-even)\\
Boundary levels $10,14$ do not win
  & \VerifiedR{$n\le36$ exact}\\
$\dn^3B_n\in\Z$ for cuspidal companions
  & \VerifiedR{$n\le36$--$40$} --- \emph{not} proved\\
$e_p=3h_p$ at generic $p$, i.e.\ $\kappa=3$
  & \VerifiedR{$n\le72$, both $F$}\\
Universal $\dn^2$-type theorem
  & \ExcludedR{generic primes}\\
$v_5(B_{5m})=v_5(B_m)-2$
  & \VerifiedR{$m\le11$} --- \Open{} as S-3\\
Integral-model extremal problem
  & \Open\\
Catalan / Koutschan--Zudilin outside the modular class
  & \ExcludedR{bounded scans}\\
\bottomrule
\end{tabular}
\medskip
\caption{Evidence ledger for the claims of this paper.}
\label{tab:evidence}
\end{table}

\section{Honest limits}
\label{sec:limits}

This section exists because the central result is a negative, and a
negative is only as good as its stated domain. Everything below is a
limitation of the present work, not a hedge.

\subsection{Search bounds}

\begin{enumerate}[label=(L\arabic*),leftmargin=3em]
\item \textbf{Levels $\le25$.} Larger levels were not scanned. There is
  no theoretical bar at $25$; it is where genus-zero standard eta
  hauptmoduln run out in the tabulated list used.
\item \textbf{One hauptmodul per level.} Exactly one standard integral
  eta-quotient coordinate was scanned per level --- with the sole
  exception of level $5$, whose two normalisations straddle the
  win/no-gain line. That exception is not a nuisance: it is the
  \emph{proof} that the single-coordinate scan is incomplete in principle
  (Lesson~\ref{lesson:norm}), and it is why
  Problem~\ref{prob:extremal} is open rather than closed.
\item \textbf{Seeded zero-finding.} The Newton search for critical values
  is seeded. At levels $6$, $12$ and $18$ the second class was taken from
  the known family operators or from scale estimates rather than found
  fresh. For levels $6$ and $12$ these values are independently known and
  exact; for level $18$ only the \emph{scale} is used, which suffices to
  rule out a win but does not resolve the second class.
\item \textbf{Odd-part dimensions not split in general.} The scan reports
  $\dim S_4(\Ga(N))$, not $\dim S_4^\pm$. The splitting was carried out
  only where it mattered --- level $6$ (where $\dim S_4=1$ and the sign is
  $+$) and level $10$. At other losing levels the odd dimension is
  irrelevant, since they lose the race regardless.
\item \textbf{Weight $r=2$ under-scanned.} The weight-$3$-with-character
  spaces relevant to $r=2$ were not systematically scanned; the reported
  $r=2$ column is a by-product of the same fold data.
\item \textbf{Level 14's coordinate.} Mis-specified in the first run and
  corrected afterwards (see \S\ref{sec:s1}). The corrected level sits at
  the same unit-circle scale, so the verdict is unchanged --- but the
  incident is a reminder that hauptmodul bookkeeping is error-prone and
  every entry of Table~\ref{tab:scan} should be re-derived by any reader
  who wants to build on it.
\end{enumerate}

\subsection{What is assumed rather than proved}

\begin{enumerate}[label=(A\arabic*),leftmargin=3em]
\item \textbf{The $\dn^r$ cost.} The identification $\delta=e^r$ assumes
  that the cuspidal companion's denominators really are governed by
  $\dn^r$. For the Eisenstein (classical) cases this is proved. For the
  cuspidal companions it is \emph{observed}: exactly to $n\le40$ at
  levels $12$ and $6$, to $n\le36$ at level $10$, and prime-locally to
  $n\le72$. It is not a theorem. A reader looking for a way through
  Problem~\ref{prob:cell} should note that this cuts both ways: the
  cuspidal denominators could in principle be \emph{worse} than $\dn^r$
  at some level, which would only strengthen the negative --- or better at
  some unexamined level, which would weaken it.
\item \textbf{The rate dictionary.} The identifications
  $\alpha=1/|\tc|$, $|\beta|=1/|\tcp|$ are the standard singularity
  analysis of the Picard--Fuchs operator, valid when the source vanishes
  at the fold to the required order and the fold is a regular singular
  point with the expected exponents. They were checked against the two
  classical constructions and the manufactured level-$12$ apparatus. They
  are not re-derived here for each scanned level.
\item \textbf{Numerical fold data.} Table~\ref{tab:scan}'s entries are
  numerical, at $50$ digits working precision, with a per-level M\"obius
  residual check. The exact values are known independently only at levels
  $5$, $6$, $8$, $12$ and $20$.
\item \textbf{$\Phi_\gamma(\tau_c)=0$.} The numerical statement is
  overwhelming ($10^{-122}$), and the parity argument that explains it is
  a proof \emph{given} the Atkin--Lehner eigenvalue of $\Phi_\gamma$,
  which is read off the Eisenstein decomposition rather than derived here
  from first principles. We grade the combination \Proved{} but flag the
  input.
\item \textbf{Mechanism claims.} The two effects of fold-vanishing
  (\S\ref{sec:unification}(i),(ii)) are argued and observed; the
  period-parity half in particular has a written route with one
  normalisation constant still open in the source ledger. Nothing in this
  paper's verdict depends on that constant, since the verdict is about
  rates.
\end{enumerate}

\subsection{What the negative does and does not say}

It says: within an explicit finite family of modular constructions ---
one standard integral coordinate per genus-zero level $N\le25$, odd
sources, Atkin--Lehner folds --- Ap\'ery's level-$6$ Eisenstein
construction is the only one whose score is below $1$, and no level with
a nonzero odd cusp space comes within a factor of $4$.

It does not say that no cuspidal irrationality proof exists, nor even
that none exists by this mechanism. Two doors are explicitly open:
Problem~\ref{prob:extremal} (a different integral coordinate, at any
level, possibly large), and the possibility of a structurally different
mechanism altogether --- for which the two hypergeometric control
experiments of \S\ref{sec:unification} suggest the known candidates in
the literature are not it.

\section{Closing: what a scan is for}
\label{sec:closing}

A scan of this kind produces three kinds of output, and it is worth
separating them.

\emph{A verdict.} The cuspidal-Ap\'ery cell is empty in the scanned
domain, analytically and arithmetically. If one is looking for a second
Ap\'ery, one should not look here.

\emph{A reframing.} Before the scan, ``why is Ap\'ery's source
Eisenstein?'' had no answer beyond ``that is what the computation gives''.
After Finding~\ref{find:phi} and Observation~\ref{obs:unify}, it has one:
Ap\'ery's construction is the odd-source-vanishing-at-the-fold mechanism,
and at the only level where that mechanism wins its race, the odd part of
the weight-$4$ space is purely Eisenstein. The Eisenstein-ness is forced.
Table~\ref{tab:2x2} is the whole story, and three of its four cells are
now understood.

\emph{New arithmetic that nobody was looking for.} The scan was a search
for a race winner and found none; but en route it produced the third
instance of the oldspace vanishing lemma, the $\kappa=3$ generic
denominator law, the exclusion of the universal $\dn^2$ theorem, and the
exact $5$-adic descent law $v_5(B_{5m})=v_5(B_m)-2$ --- an equality, not
an inequality, with two identified proof routes. In this programme
bounded negatives are deliverables, and the structure they force out on
the way is usually the durable part.

What remains, precisely: prove the descent law (Problem~\ref{prob:s3}),
and solve the extremal problem over integral models
(Problem~\ref{prob:extremal}). Failing both, a cuspidal irrationality
proof needs a mechanism this paper has not scanned.

\subsection*{Acknowledgements}

The score formulation of the race, the scanner proposal, the protocol for
the prime-local denominator split, and the extremal-function framing of
the surviving open problem are all due to our collaborator Sol.

\appendix

\section{Provenance}
\label{app:prov}

For a reader wishing to reproduce or extend the computations, the
underlying records of this programme are:

\begin{itemize}[leftmargin=1.4em]
\item \texttt{work/SCANNER\_LEDGER.md} --- the scan itself: method
  (\S1), the table (\S2), the verdict (\S3), the normalisation lesson
  (\S4), honest limits (\S5), the level-$10$ boundary test (\S6), the
  prime-local law (\S7), and the $5$-adic descent law (\S8). All of
  \S\ref{sec:scan}--\S\ref{sec:limits} above tracks this ledger.
\item \texttt{work/D1\_CUSPIDAL\_APPARATUS.md} --- the manufactured
  cuspidal apparatus, and (\S9) the race framework together with the
  $\Phi_\gamma(\tau_c)=0$ discovery and the $2\times2$ table.
\item \texttt{work/CATALAN\_CONTROL.md},
  \texttt{work/KZ\_ELLIPTIC\_REPORT.md} --- the hypergeometric boundary:
  why the modular class is the arena.
\item \texttt{work/z5eps/eps65\_scanner.py} --- the scanner: eta
  hauptmodul table, numerical Atkin--Lehner M\"obius fit, fixed points,
  $\dim S_k(\Ga(N))$, and the score columns for $r=2,3$. The remaining
  computations (levels $10$, $14$; the prime-local splits; the descent
  law probe) were run inline in the same session's exact-arithmetic
  environment.
\end{itemize}

The companion papers of this arc are the cuspidal-companion paper (the
$L(f_6,3)/2$ apparatus, the oldspace lemma, the parity mechanism and the
six-item proof package) and the modular-anchor papers of the wider
programme, to which we refer for the theory of sources
$\Phi=t\sigma^r/(PF)$ and the fold connection formula.

\begin{thebibliography}{9}

\bibitem{Apery} R. Ap\'ery, \emph{Irrationalit\'e de $\zeta(2)$ et
$\zeta(3)$}, Ast\'erisque \textbf{61} (1979), 11--13.

\bibitem{vdP} A. van der Poorten, \emph{A proof that Euler missed\dots\
Ap\'ery's proof of the irrationality of $\zeta(3)$}, Math. Intelligencer
\textbf{1} (1979), 195--203.

\bibitem{Beukers} F. Beukers, \emph{Irrationality proofs using modular
forms}, Ast\'erisque \textbf{147--148} (1987), 271--283.

\bibitem{Zagier} D. Zagier, \emph{Integral solutions of Ap\'ery-like
recurrence equations}, in: Groups and Symmetries, CRM Proc. Lecture
Notes \textbf{47} (2009), 349--366.

\bibitem{ZudilinCatalan} W. Zudilin, \emph{An Ap\'ery-like difference
equation for Catalan's constant}, Electron. J. Combin. \textbf{10}
(2003), \#R14; arXiv:math/0201024.

\bibitem{KZ} C. Koutschan and W. Zudilin, \emph{Ap\'ery limits for
elliptic $L$-values}, Bull. Aust. Math. Soc. \textbf{106} (2022),
273--279; arXiv:2111.08796.

\bibitem{AtkinLehner} A. O. L. Atkin and J. Lehner, \emph{Hecke operators
on $\Gamma_0(m)$}, Math. Ann. \textbf{185} (1970), 134--160.

\bibitem{D1} [author] and collaborators, \emph{A cuspidal companion: an
Ap\'ery-quality apparatus for a critical $L$-value}, companion paper of
this arc; technical record \texttt{work/D1\_CUSPIDAL\_APPARATUS.md}.

\end{thebibliography}

\end{document}
